\documentclass[11pt,a4paper]{article}
\usepackage[utf8]{inputenc}
\usepackage[T1]{fontenc}
\usepackage{geometry}
\usepackage{listings}
\usepackage{hyperref}
\usepackage{xcolor}
\usepackage{array}
\usepackage{booktabs}
\usepackage{amsmath}
\usepackage{pdflscape}
\geometry{margin=1in}

\title{Java Subset Lexical Analyzer\\\large Design and Implementation Report}
\author{\textbf{Participants:}\n
Braa Mostafa (22010636)\n
Mohamed Sabry (22011129)\n
Muhnad Sayed (22011902)\n
Youssef Khamis (22011381)}
\date{\today}

\begin{document}
\maketitle
\newpage

\section{Project Overview}
This project implements a complete, table-driven lexical analyzer for a subset of
Java, written in C++17. The program reads a \emph{rule file} that describes the
token classes of the language with regular expressions, regular definitions,
keyword groups, and punctuation groups, and from this description it
automatically constructs a minimal deterministic finite automaton (DFA). The
resulting DFA is then used to scan Java-subset source programs and emit the
stream of recognized token classes.

The implemented system follows the classic scanner-generator pipeline described
in the Dragon Book:
\begin{enumerate}
    \item \textbf{Rule parsing} -- read and validate the lexical rules, expand
          named regular definitions, and assign priorities (keywords first).
    \item \textbf{Regex parsing} -- convert each regular expression into an
          abstract syntax tree (AST) with a recursive-descent parser.
    \item \textbf{NFA construction} -- translate each AST into an NFA fragment
          using Thompson's construction and combine all fragments with a common
          start state.
    \item \textbf{Subset construction} -- determinize the combined NFA into a
          DFA, resolving accepting-rule conflicts by priority.
    \item \textbf{DFA minimization} -- merge equivalent states with partition
          refinement (Moore's algorithm) to obtain the minimal DFA.
    \item \textbf{Scanning} -- run the maximal-munch (longest-match) algorithm
          over the input and emit one token class per recognized lexeme.
\end{enumerate}

For the assignment's Phase~1 rule file, the pipeline produces concrete,
verifiable numbers: the $18$ token rules are compiled into a Thompson NFA with
$153$ states, subset construction yields a DFA with $54$ states, and
minimization reduces it to the \textbf{minimal DFA with 44 states} over an
alphabet of $77$ input characters. The transition table of this minimal DFA is
written to the output file required by the assignment, and the token stream
produced for the example test program matches the expected output exactly.

\section{Project Structure}
The lexer lives under \texttt{src/lexer/} and is split into small, independently
testable components:
\begin{itemize}
    \item \texttt{rules.hpp} / \texttt{rules.cpp}: the \texttt{Rule} structure and
          the \texttt{RuleReader} that parses the rule file, expands named
          definitions, and assigns priorities.
    \item \texttt{automata.hpp}: the shared automata data structures
          (\texttt{NfaState}, \texttt{Nfa}, \texttt{DfaState}, \texttt{Dfa}) and the
          builder interfaces.
    \item \texttt{nfa.cpp}: regular-expression tokenizer, recursive-descent regex
          parser, and Thompson NFA construction (\texttt{NfaBuilder}).
    \item \texttt{dfa.cpp}: epsilon-closure computation, subset construction,
          partition-refinement minimization, and transition-table printing
          (\texttt{DfaBuilder}).
    \item \texttt{scanner.hpp} / \texttt{scanner.cpp}: the \texttt{Scanner} class
          implementing maximal-munch scanning, the identifier symbol table, and
          panic-mode error recovery.
    \item \texttt{lexer.hpp} / \texttt{lexer.cpp}: the \texttt{LexerDefinition}
          facade that builds the rules and the minimized DFA once per process and
          hands out a fresh \texttt{Scanner} per source stream.
    \item \texttt{main.cpp}: the command-line driver
          (\texttt{lexer RULES\_FILE TABLE\_FILE SOURCE\_FILE...}).
    \item \texttt{test\_files/}: the Phase~1 rule file, the example test program,
          and its expected token stream.
    \item \texttt{tests/}: unit and end-to-end tests (\texttt{test\_nfa},
          \texttt{test\_dfa}, \texttt{test\_scanner}, \texttt{test\_e2e},
          \texttt{test\_flex\_parity}).
    \item \texttt{bonus\_flex/lexer.l}: the Flex specification for the bonus task,
          kept behaviorally identical to the hand-written lexer.
\end{itemize}

\section{Description of the Used Data Structures}

\subsection{Rule}
\begin{lstlisting}[basicstyle=\ttfamily\small,language=C++]
struct Rule {
    std::string name;        // token class, e.g. "id", "while"
    std::string expression;  // fully expanded regular expression
    int priority;            // smaller wins; rule index breaks ties
};
\end{lstlisting}
A \texttt{Rule} carries the token-class name, the regular expression \emph{after}
all named definitions have been expanded, and an integer priority. The rule
reader reserves the lowest (strongest) priorities for keywords, so reserved
words such as \texttt{int} and \texttt{while} always beat the generic
\texttt{id} pattern. Rule vector indices stay stable through every later stage;
the DFA stores these indices in its accepting states.

\subsection{NFA structures}
\begin{lstlisting}[basicstyle=\ttfamily\small,language=C++]
struct NfaState {
    std::map<char, std::vector<int>> transitions; // labeled edges
    std::vector<int> epsilon;                     // epsilon edges
    int rule = -1;                                // accepting rule or -1
};
struct Nfa { std::vector<NfaState> states; int start = 0; };
\end{lstlisting}
States are identified by their index in the \texttt{states} vector, so an edge is
just an \texttt{int}. Each state keeps a character-indexed map of possible target
lists (a state may move to several states on the same character --
nondeterminism) plus a separate list of $\varepsilon$-moves. The one accepting
state of each Thompson fragment stores the index of the rule it accepts.

\subsection{Regex AST (internal to the NFA builder)}
The regex compiler first tokenizes an expression into \texttt{Token} records
(literal, escaped literal, \texttt{\textbackslash L} for $\varepsilon$, parentheses,
\texttt{|}, \texttt{*}, \texttt{+}, and the range dash), then builds a tree of
\texttt{AstNode} objects with kinds \texttt{Literal}, \texttt{Epsilon},
\texttt{Range}, \texttt{Concat}, \texttt{Alt}, \texttt{Star}, and \texttt{Plus}.
Nodes own their children through \texttt{std::unique\_ptr}, so the tree frees
itself. A small \texttt{Fragment \{int start; int accept;\}} pair is used while
translating the AST into NFA states.

\subsection{DFA structures}
\begin{lstlisting}[basicstyle=\ttfamily\small,language=C++]
struct DfaState {
    std::map<char, int> transitions; // missing edge = failure
    int rule = -1;                   // winning rule or -1
};
struct Dfa { std::vector<DfaState> states; int start = 0; };
\end{lstlisting}
A DFA state has at most one target per character, so the map value is a single
\texttt{int}. A character absent from the map means the automaton dies on that
input -- the dead state is represented implicitly rather than materialized
(printed as `\texttt{-}' in the transition table). The \texttt{rule} field holds
the highest-priority rule among the NFA accepting states merged into this DFA
state.

\subsection{Containers used by the algorithms}
\begin{itemize}
    \item \texttt{std::set<int>}: subsets of NFA states ($\varepsilon$-closures).
          Sets are ordered and canonical, so equal subsets compare equal.
    \item \texttt{std::map<std::set<int>, int>}: maps each discovered NFA-state
          subset to its DFA state id during subset construction, guaranteeing
          each subset becomes exactly one DFA state.
    \item \texttt{std::queue<...>}: the worklist for the breadth-first
          $\varepsilon$-closure computation and for subset construction.
    \item \texttt{std::vector<std::set<int>>} and
          \texttt{std::map<std::map<char,int>, std::set<int>>}: the partitions and
          the signature-based splitting map used by the minimizer.
    \item \texttt{std::map<std::string, Definition>} plus memoization maps inside
          the rule reader for named-definition expansion and cycle detection.
\end{itemize}

\subsection{Scanner}
\begin{lstlisting}[basicstyle=\ttfamily\small,language=C++]
struct ScannedToken {
    std::string type;    // token class name
    std::string lexeme;  // matched text
    int symbol;          // identifier's symbol-table index, or -1
    std::size_t offset;  // zero-based byte offset in the input
};
\end{lstlisting}
The \texttt{Scanner} walks the DFA over an \texttt{std::istream}. It keeps:
\begin{itemize}
    \item a \texttt{std::deque<char> pending\_} push-back buffer: characters read
          past the last accepting state are ``unread'' into this deque so the
          next token starts exactly after the previous lexeme (required by
          maximal munch);
    \item a \texttt{std::vector<std::string> symbols\_} acting as a simple symbol
          table -- every distinct identifier lexeme is interned once and its
          index is attached to the token;
    \item a \texttt{std::vector<std::size\_t> error\_offsets\_} recording the byte
          offset of every lexical error for diagnostics.
\end{itemize}
Finally, \texttt{LexerDefinition} bundles the rule vector and the minimized DFA
so they are built once and reused for any number of source files.

\section{Explanation of All Algorithms and Techniques Used}

\subsection{Rule-file parsing and definition expansion}
The rule file is read line by line. Four line forms are supported, following the
handout:
\begin{itemize}
    \item \texttt{name = expr} -- a \emph{regular definition} (e.g.
          \texttt{letter = a-z | A-Z}); it does not create a token by itself.
    \item \texttt{name : expr} -- a \emph{token rule} (e.g.
          \texttt{id: letter (letter|digit)*}).
    \item \texttt{\{w1 w2 ...\}} -- a \emph{keyword group}; every word becomes a
          reserved-word token whose pattern is the word itself.
    \item \texttt{[p1 p2 ...]} -- a \emph{punctuation group}; every
          (possibly escaped) symbol becomes a one-character token.
\end{itemize}
Group members are turned into safe regular expressions by escaping every regex
metacharacter (\texttt{escape\_literal\_expression}). Identifier-like words of
two or more characters inside an expression are treated as references to named
definitions and are expanded \emph{recursively with memoization}; the expansion
is wrapped in parentheses to preserve operator precedence. Range endpoints such
as \texttt{a} in \texttt{a-z} are deliberately \emph{not} expanded even if a
definition with the same single-letter name exists. The expander detects
\textbf{undefined references}, \textbf{duplicate definitions}, and
\textbf{cyclic definitions} and reports them as errors with the offending line
number. Priorities are then assigned: the $k$ keywords receive priorities
$0..k-1$ and all other rules continue from $k$, in file order, so reserved words
always dominate.

\subsection{Regular-expression parsing (recursive descent)}
Each expanded expression is tokenized (whitespace between symbols is ignored;
\texttt{\textbackslash x} yields the literal \texttt{x}; \texttt{\textbackslash L}
yields $\varepsilon$) and parsed with a standard precedence-layered
recursive-descent grammar:
\[
\textit{alt} \rightarrow \textit{concat}\ (\texttt{|}\ \textit{concat})^* \qquad
\textit{concat} \rightarrow \textit{postfix}^+ \qquad
\textit{postfix} \rightarrow \textit{atom}\ (\texttt{*}\,|\,\texttt{+})^* \qquad
\textit{atom} \rightarrow \texttt{(}\textit{alt}\texttt{)}\ |\ \text{literal}\ |\ \text{range}\ |\ \varepsilon
\]
This gives the conventional precedence (closure $>$ concatenation $>$ union).
Character ranges \texttt{lo-hi} are recognized at the atom level and validated
(\texttt{lo} $\le$ \texttt{hi}). Malformed inputs -- unmatched parentheses,
dangling \texttt{*}/\texttt{+}, empty alternatives, trailing backslashes --
raise descriptive errors naming the offending rule.

\subsection{Thompson construction (regex $\rightarrow$ NFA)}
The AST is translated bottom-up into NFA fragments with the textbook Thompson
rules: a literal creates two states joined by one labeled edge; a range
\texttt{lo-hi} creates one edge per character in the range; concatenation links
fragment~$A$'s accept to fragment~$B$'s start with an $\varepsilon$-edge;
union, star, and plus add new start/accept states wired with
$\varepsilon$-edges in the usual way (\texttt{+} is like \texttt{*} but without
the ``skip'' edge, so at least one iteration is required). Every rule's fragment
accept state is tagged with the rule index, and finally a single new
\textbf{global start state} is created with $\varepsilon$-edges to every
fragment start, producing one combined NFA that matches any token.

As a safety check, after building each fragment the builder runs a BFS over
$\varepsilon$-edges (\texttt{can\_reach\_eps}) from the fragment start to its
accept state; if the accept state is $\varepsilon$-reachable the expression can
match the empty string, which would make the scanner loop forever, so such rules
are rejected at construction time.

\subsection{Subset construction (NFA $\rightarrow$ DFA)}
Determinization uses the standard worklist algorithm:
\begin{enumerate}
    \item The DFA start state is $\varepsilon\text{-closure}(\{s_0\})$, computed by a
          BFS over $\varepsilon$-edges with a \texttt{std::queue}.
    \item For each subset $T$ popped from the worklist and each character $c$
          that actually labels an edge out of $T$ (the local alphabet, not all
          256 bytes), compute
          $\varepsilon\text{-closure}(\mathit{move}(T, c))$.
    \item If the resulting subset is new, register it in
          \texttt{std::map<std::set<int>,int>} and push it on the worklist; in
          all cases record the DFA transition.
    \item The \emph{empty} subset is never materialized: a transition into it is
          simply omitted, which implicitly encodes the dead state.
\end{enumerate}
For Phase~1 this turns the 153-state NFA into a 54-state DFA.

\subsection{Accepting-rule conflict resolution}
A DFA state may contain several accepting NFA states (e.g.\ after reading
\texttt{int} the subset contains both the \texttt{int}-keyword accept and an
\texttt{id} accept). While each subset is processed, the builder scans its NFA
states and keeps the rule with the \textbf{smallest priority}; if two rules have
equal priority the smaller original rule index wins, making the result fully
deterministic. Because keywords received the smallest priorities during rule
parsing, this resolves the classic keyword-vs-identifier ambiguity in favor of
the reserved words -- inside the automaton itself, with no special cases in the
scanner.

\subsection{DFA minimization (partition refinement / Moore's algorithm)}
Minimization starts from the initial partition that groups states \emph{by their
accepting rule id} (one block per token class plus one block for
non-accepting states). Grouping by rule rather than by a mere
accepting/non-accepting flag is essential: merging an \texttt{int}-accepting
state with a \texttt{while}-accepting state would be language-preserving for the
union automaton but would destroy the token classification. The algorithm then
repeats:
\begin{enumerate}
    \item map every state to the index of its current block;
    \item compute for every state a \emph{signature}
          $c \mapsto \text{block}(\delta(s,c))$ over the global alphabet, with a
          missing transition encoded as block $-1$;
    \item split each block by signature equality
          (\texttt{std::map<std::map<char,int>, std::set<int>>});
\end{enumerate}
until no block splits. Each final block becomes one state of the minimal DFA;
transitions are remapped through a representative, the block containing the old
start state becomes the new start, and the accepting rule id of the block is
preserved. For Phase~1 this reduces 54 states to the minimal \textbf{44 states}.

\subsection{Maximal-munch (longest-match) scanning}
\texttt{Scanner::nextToken()} implements the classical two-pointer maximal-munch
algorithm:
\begin{enumerate}
    \item skip whitespace (space, tab, CR, LF, form feed, vertical tab), counting
          offsets; return \texttt{std::nullopt} on end of input;
    \item walk the DFA from the start state, appending every examined character
          to a buffer; each time an accepting state is entered, remember the rule
          and the current length (\emph{last-accept bookkeeping});
    \item stop at the first missing transition or at end of input, then
          \emph{unread} every character after the last accepted length back into
          the \texttt{pending\_} deque so it will be rescanned as the start of
          the next token;
    \item emit the token: its class name, its lexeme, its byte offset, and -- for
          \texttt{id} tokens -- its index in the interned symbol table.
\end{enumerate}
This guarantees, for example, that \texttt{!=10} is split into \texttt{relop}
(\texttt{!=}) followed by \texttt{num} (\texttt{10}), and that
\texttt{passX} is a single identifier rather than the keyword-prefix
\texttt{pass} plus garbage.

\subsection{Error handling (panic mode)}
If the walk fails before any accepting state was reached, the scanner records
the current byte offset in \texttt{error\_offsets\_}, \textbf{discards exactly one
byte} (guaranteeing progress and termination), pushes the rest of the examined
characters back, and resumes scanning. The driver prints every error offset to
\texttt{stderr} after the token stream, so a single bad character never aborts
the whole scan. The constructor additionally validates the DFA defensively: it
rejects an automaton whose start state accepts (empty token), out-of-range
transition targets, and invalid rule ids.

\subsection{Overall driver}
\texttt{LexerDefinition} runs
\texttt{RuleReader} $\rightarrow$ \texttt{NfaBuilder} $\rightarrow$
\texttt{DfaBuilder::build} $\rightarrow$ \texttt{DfaBuilder::minimize} once, and
\texttt{main} then writes the transition table to the requested file and scans
each source file given on the command line, printing one token class per line.

\section{Resultant Transition Table for the Minimal DFA}

\subsection{How it was produced}
Built and run from the repository root:
\begin{lstlisting}[basicstyle=\ttfamily\small]
cmake -S . -B build && cmake --build build
./build/lexer src/lexer/test_files/phase1.rules build/phase1.table \
    src/lexer/test_files/phase1_program.java > build/phase1.tokens
less -s build/phase1.table
\end{lstlisting}

\subsection{Summary of the minimal DFA}
For the Phase~1 rules the pipeline produces:
\begin{center}
\begin{tabular}{lr}
\toprule
Token rules & 18 \\
Thompson NFA states & 153 \\
DFA states after subset construction & 54 \\
\textbf{States in the minimal DFA} & \textbf{44} \\
Alphabet size (characters with transitions) & 77 \\
Start state & 3 (marked \texttt{->} in the table) \\
\bottomrule
\end{tabular}
\end{center}

\subsection{Table format and rule numbering}
Rows are DFA states ($0$--$43$); columns are the input characters; a cell holds
the target state or `\texttt{-}' for a missing edge (implicit dead state). The
last column, \texttt{Rule}, is the id of the accepted token rule, or
\texttt{None} for non-accepting states. Rule ids map to token classes as
follows (priority in parentheses; keywords hold the strongest priorities):
\begin{center}
\small
\begin{tabular}{lll@{\hspace{2em}}lll}
\toprule
Id & Token & Priority & Id & Token & Priority \\
\midrule
0 & \texttt{id}      & 6 & 9  & \texttt{while}  & 5 \\
1 & \texttt{boolean} & 0 & 10 & \texttt{;}      & 10 \\
2 & \texttt{int}     & 1 & 11 & \texttt{,}      & 11 \\
3 & \texttt{float}   & 2 & 12 & \texttt{(}      & 12 \\
4 & \texttt{num}     & 7 & 13 & \texttt{)}      & 13 \\
5 & \texttt{relop}   & 8 & 14 & \texttt{\{}     & 14 \\
6 & \texttt{assign}  & 9 & 15 & \texttt{\}}     & 15 \\
7 & \texttt{if}      & 3 & 16 & \texttt{addop}  & 16 \\
8 & \texttt{else}    & 4 & 17 & \texttt{mulop}  & 17 \\
\bottomrule
\end{tabular}
\end{center}

The table itself illustrates the techniques described above. For instance,
state~3 is the start state; reading \texttt{i} leads to state~11, then
\texttt{n} to state~14, then \texttt{t} to state~22, which accepts rule~2
(\texttt{int}) -- yet any further letter or digit moves to state~4, the generic
identifier state accepting rule~0 (\texttt{id}). That is exactly the
keyword-versus-identifier resolution and the longest-match behavior encoded in
the automaton. Likewise state~0 (reached from state~3 on \texttt{!}) only has an
edge on \texttt{=}, so a lone \texttt{!} is a lexical error while \texttt{!=}
accepts \texttt{relop}.

\subsection{The exact generated table}
The following is the verbatim content of the generated table file
(as displayed by \texttt{less -s build/phase1.table}). To improve readability,
we present it in landscape format and with a compact monospaced style:

\begin{landscape}
\small
\begin{verbatim}
State    !   (   )   *   +   ,   -   .   /   0   1   2   3   4   5   6   7   8   9   ;   <   =   >   A   B   C   D   E   F   G   H   I   J   K   L   M   N   O   P   Q   R   S   T   U   V   W   X   Y   Z   a   b   c   d   e   f   g   h   i   j   k   l   m   n   o   p   q   r   s   t   u   v   w   x   y   z   {   }   Rule
-----------------------------------------------------------------------------------------------------------------------------------------------------------------------------------------------------------------------------------------------------------------------------------------------------------------------------------
0        -   -   -   -   -   -   -   -   -   -   -   -   -   -   -   -   -   -   -   -   -  30   -   -   -   -   -   -   -   -   -   -   -   -   -   -   -   -   -   -   -   -   -   -   -   -   -   -   -   -   -   -   -   -   -   -   -   -   -   -   -   -   -   -   -   -   -   -   -   -   -   -   -   -   -   -  None
1        -   -   -   -   -   -   -   -   -  27  27  27  27  27  27  27  27  27  27   -   -   -   -   -   -   -   -   -   -   -   -   -   -   -   -   -   -   -   -   -   -   -   -   -   -   -   -   -   -   -   -   -   -   -   -   -   -   -   -   -   -   -   -   -   -   -   -   -   -   -   -   -  None
2        -   -   -   -   -   -   -   -   -  28  28  28  28  28  28  28  28  28  28   -   -   -   -   -   -   -   -   -   -   -   -   -   -   -   -   -   -   -   -   -   -   -   -   -   -   -   -   -   -   -   -   -   -   -   -   -   -   -   -   -   -   -   -   -   -   -   -   -   -   -   -   -  None
->3       0  38  39  43  42  37  42  -  43  29  29  29  29  29  29  29  29  29  29  36  31  32  31   4   4   4   4   4   4   4   4   4   4   4   4   4   4   4   4   4   4   4   4   4   4   4   4   4   4   4   5   4   4  11   7   4   4  21   4   4   4   4   4   4   4   4   4   4   4   4   4   9   4   4   4  40  41  None
4        -   -   -   -   -   -   -   -   -   4   4   4   4   4   4   4   4   4   4   -   -   -   -   4   4   4   4   4   4   4   4   4   4   4   4   4   4   4   4   4   4   4   4   4   4   4   4   4   4   4   4   4   4   4   4   4   4   4   4   4   4   4   4   4   4   4   4   4   4   4   4   -   -   0
5        -   -   -   -   -   -   -   -   -   4   4   4   4   4   4   4   4   4   4   -   -   -   -   4   4   4   4   4   4   4   4   4   4   4   4   4   4   4   4   4   4   4   4   4   4   4   4   4   4   4   4   4   4   4   4   4   4   4   4   4   4   4   4   4   6   4   4   4   4   4   4   4   4   4   4   4   4   -   -   0
6        -   -   -   -   -   -   -   -   -   4   4   4   4   4   4   4   4   4   4   -   -   -   -   4   4   4   4   4   4   4   4   4   4   4   4   4   4   4   4   4   4   4   4   4   4   4   4   4   4   4   4   4   4   4   4   4   4   4   4   4   4   4   4   4   8   4   4   4   4   4   4   4   4   4   4   4   4   -   -   0
7        -   -   -   -   -   -   -   -   -   4   4   4   4   4   4   4   4   4   4   -   -   -   -   4   4   4   4   4   4   4   4   4   4   4   4   4   4   4   4   4   4   4   4   4   4   4   4   4   4   4   4   4   4   4   4   4   4   4   4  10   4   4   4   4   4   4   4   4   4   4   4   4   4   4   4   4   -   -   0
8        -   -   -   -   -   -   -   -   -   4   4   4   4   4   4   4   4   4   4   -   -   -   -   4   4   4   4   4   4   4   4   4   4   4   4   4   4   4   4   4   4   4   4   4   4   4   4   4   4   4   4   4   4   4   4   4   4   4   4   4   4   4  13   4   4   4   4   4   4   4   4   4   4   4   4   4   -   -   0
9        -   -   -   -   -   -   -   -   -   4   4   4   4   4   4   4   4   4   4   -   -   -   -   4   4   4   4   4   4   4   4   4   4   4   4   4   4   4   4   4   4   4   4   4   4   4   4   4   4   4   4   4   4   4   4   4   4   4  12   4   4   4   4   4   4   4   4   4   4   4   4   4   4   4   4   -   -   0
10       -   -   -   -   -   -   -   -   -   4   4   4   4   4   4   4   4   4   4   -   -   -   -   4   4   4   4   4   4   4   4   4   4   4   4   4   4   4   4   4   4   4   4   4   4   4   4   4   4   4   4   4   4   4   4   4   4   4   4   4   4   4   4   4  16   4   4   4   4   4   4   4   4   4   4   4   4   -   -   0
11       -   -   -   -   -   -   -   -   -   4   4   4   4   4   4   4   4   4   4   -   -   -   -   4   4   4   4   4   4   4   4   4   4   4   4   4   4   4   4   4   4   4   4   4   4   4   4   4   4   4   4   4   4   4   4   4   4   4   4   4   4   4   4  14   4   4   4   4   4   4   4   4   4   4   4   4   -   -   0
12       -   -   -   -   -   -   -   -   -   4   4   4   4   4   4   4   4   4   4   -   -   -   -   4   4   4   4   4   4   4   4   4   4   4   4   4   4   4   4   4   4   4   4   4   4   4   4   4   4   4   4   4   4   4   4   4   4   4   4   4   4  15   4   4   4   4   4   4   4   4   4   4   4   4   4   -   -   0
13       -   -   -   -   -   -   -   -   -   4   4   4   4   4   4   4   4   4   4   -   -   -   -   4   4   4   4   4   4   4   4   4   4   4   4   4   4   4   4   4   4   4   4   4   4   4   4   4   4   4   4   4   4   4  17   4   4   4   4   4   4   4   4   4   4   4   4   4   4   4   4   4   4   4   4   -   -   0
14       -   -   -   -   -   -   -   -   -   4   4   4   4   4   4   4   4   4   4   -   -   -   -   4   4   4   4   4   4   4   4   4   4   4   4   4   4   4   4   4   4   4   4   4   4   4   4   4   4   4   4   4   4   4   4   4   4   4   4   4   4   4   4   4   4   4   4  22   4   4   4   4   4   4   4   -   -   0
15       -   -   -   -   -   -   -   -   -   4   4   4   4   4   4   4   4   4   4   -   -   -   -   4   4   4   4   4   4   4   4   4   4   4   4   4   4   4   4   4   4   4   4   4   4   4   4   4   4   4   4   4   4   4   4   4   4   4   4   4   4  23   4   4   4   4   4   4   4   4   4   4   4   4   4   -   -   0
16       -   -   -   -   -   -   -   -   -   4   4   4   4   4   4   4   4   4   4   -   -   -   -   4   4   4   4   4   4   4   4   4   4   4   4   4   4   4   4   4   4   4   4   4   4   4   4   4  19   4   4   4   4   4   4   4   4   4   4   4   4   4   4   4   4   4   4   4   4   4   4   4   4   4   4   -   -   0
17       -   -   -   -   -   -   -   -   -   4   4   4   4   4   4   4   4   4   4   -   -   -   -   4   4   4   4   4   4   4   4   4   4   4   4   4   4   4   4   4   4   4   4   4   4   4   4   4  20   4   4   4   4   4   4   4   4   4   4   4   4   4   4   4   4   4   4   4   4   4   4   4   4   4   4   -   -   0
18       -   -   -   -   -   -   -   -   -   4   4   4   4   4   4   4   4   4   4   -   -   -   -   4   4   4   4   4   4   4   4   4   4   4   4   4   4   4   4   4   4   4   4   4   4   4   4   4   4   4   4   4   4   4   4   4   4   4   4   4   4   4   4   4   4   4   4  25   4   4   4   4   4   4   4   -   -   0
19       -   -   -   -   -   -   -   -   -   4   4   4   4   4   4   4   4   4   4   -   -   -   -   4   4   4   4   4   4   4   4   4   4   4   4   4   4   4   4   4   4   4   4   4   4   4   4   4   4   4   4   4   4   4   4   4   4   4   4   4   4   4   4   4   4   4   4  26   4   4   4   4   4   4   4   -   -   0
20       -   -   -   -   -   -   -   -   -   4   4   4   4   4   4   4   4   4   4   -   -   -   -   4   4   4   4   4   4   4   4   4   4   4   4   4   4   4   4   4   4   4   4   4   4   4   4   4   4   4   4   4   4   4   4   4   4   4   4   4   4   4   4  24   4   4   4   4   4   4   4   4   4   4   4   4   -   -   0
21       -   -   -   -   -   -   -   -   -   4   4   4   4   4   4   4   4   4   4   -   -   -   -   4   4   4   4   4   4   4   4   4   4   4   4   4   4   4   4   4   4   4   4   4   4   4   4   4   4   4   4   4   4   4  33   4   4   4   4   4   4   4  18   4   4   4   4   4   4   4   4   4   4   4   4   -   -   0
22       -   -   -   -   -   -   -   -   -   4   4   4   4   4   4   4   4   4   4   -   -   -   -   4   4   4   4   4   4   4   4   4   4   4   4   4   4   4   4   4   4   4   4   4   4   4   4   4   4   4   4   4   4   4   4   4   4   4   4   4   4   4   4   4   4   4   4   4   4   4   4   -   -   0
23       -   -   -   -   -   -   -   -   -   4   4   4   4   4   4   4   4   4   4   -   -   -   -   4   4   4   4   4   4   4   4   4   4   4   4   4   4   4   4   4   4   4   4   4   4   4   4   4   4   4   4   4   4   4  35   4   4   4   4   4   4   4   4   4   4   4   4   4   4   4   4   -   -   0
24       -   -   -   -   -   -   -   -   -   4   4   4   4   4   4   4   4   4   4   -   -   -   -   4   4   4   4   4   4   4   4   4   4   4   4   4   4   4   4   4   4   4   4   4   4   4   4   4   4   4   4   4   4   4   4   4   4   4   4   4   4   4   4   4   4   4   4   4   4   4   4   -   -   1
25       -   -   -   -   -   -   -   -   -   4   4   4   4   4   4   4   4   4   4   -   -   -   -   4   4   4   4   4   4   4   4   4   4   4   4   4   4   4   4   4   4   4   4   4   4   4   4   4   4   4   4   4   4   4   4   4   4   4   4   4   4   4   4   4   4   4   4   4   4   4   4   -   -   2
26       -   -   -   -   -   -   -   -   -   4   4   4   4   4   4   4   4   4   4   -   -   -   -   4   4   4   4   4   4   4   4   4   4   4   4   4   4   4   4   4   4   4   4   4   4   4   4   4   4   4   4   4   4   4   4   4   4   4   4   4   4   4   4   4   4   4   4   4   4   4   4   -   -   3
27       -   -   -   -   -   -   -   -   -  27  27  27  27  27  27  27  27  27  27   -   -   -   -   -   -   -   -   -   -   -   -   -   -   -   -   -   -   -   -   -   -   -   -   -   -   -   -   -   -   -   -   -   -   -   -   -   -   -   -   -   -   -   -   -   -   -   -   -   -   -   -   -   4
28       -   -   -   -   -   -   -   -   -  28  28  28  28  28  28  28  28  28  28   -   -   -   -   -   -   -   -   1   -   -   -   -   -   -   -   -   -   -   -   -   -   -   -   -   -   -   -   -   -   -   -   -   -   -   -   -   -   -   -   -   -   -   -   -   -   -   -   -   -   -   -   4
29       -   -   -   -   -   -   -   2   -  29  29  29  29  29  29  29  29  29  29   -   -   -   -   -   -   -   -   -   -   -   -   -   -   -   -   -   -   -   -   -   -   -   -   -   -   -   -   -   -   -   -   -   -   -   -   -   -   -   -   -   -   -   -   -   -   -   -   -   -   -   -   4
30       -   -   -   -   -   -   -   -   -   -   -   -   -   -   -   -   -   -   -   -   -   -   -   -   -   -   -   -   -   -   -   -   -   -   -   -   -   -   -   -   -   -   -   -   -   -   -   -   -   -   -   -   -   -   -   -   -   -   -   -   -   -   -   -   -   -   -   -   -   -   -   -   -  5
31       -   -   -   -   -   -   -   -   -   -   -   -   -   -   -   -   -   -   -   -   -  30   -   -   -   -   -   -   -   -   -   -   -   -   -   -   -   -   -   -   -   -   -   -   -   -   -   -   -   -   -   -   -   -   -   -   -   -   -   -   -   -   -   -   -   -   -   -   -   -   -   -  5
32       -   -   -   -   -   -   -   -   -   -   -   -   -   -   -   -   -   -   -   -   -  30   -   -   -   -   -   -   -   -   -   -   -   -   -   -   -   -   -   -   -   -   -   -   -   -   -   -   -   -   -   -   -   -   -   -   -   -   -   -   -   -   -   -   -   -   -   -   -   -   -   -  6
33       -   -   -   -   -   -   -   -   -   4   4   4   4   4   4   4   4   4   4   -   -   -   -   4   4   4   4   4   4   4   4   4   4   4   4   4   4   4   4   4   4   4   4   4   4   4   4   4   4   4   4   4   4   4   4   4   4   4   4   4   4   4   4   4   4   4   4   4   4   4   4   -   -   7
34       -   -   -   -   -   -   -   -   -   4   4   4   4   4   4   4   4   4   4   -   -   -   -   4   4   4   4   4   4   4   4   4   4   4   4   4   4   4   4   4   4   4   4   4   4   4   4   4   4   4   4   4   4   4   4   4   4   4   4   4   4   4   4   4   4   4   4   4   4   4   4   -   -   8
35       -   -   -   -   -   -   -   -   -   4   4   4   4   4   4   4   4   4   4   -   -   -   -   4   4   4   4   4   4   4   4   4   4   4   4   4   4   4   4   4   4   4   4   4   4   4   4   4   4   4   4   4   4   4   4   4   4   4   4   4   4   4   4   4   4   4   4   4   4   4   4   -   -   9
36       -   -   -   -   -   -   -   -   -   -   -   -   -   -   -   -   -   -   -   -   -   -   -   -   -   -   -   -   -   -   -   -   -   -   -   -   -   -   -   -   -   -   -   -   -   -   -   -   -   -   -   -   -   -   -   -   -   -   -   -   -   -   -   -   -   -   -   -   -   -   -   -   -  10
37       -   -   -   -   -   -   -   -   -   -   -   -   -   -   -   -   -   -   -   -   -   -   -   -   -   -   -   -   -   -   -   -   -   -   -   -   -   -   -   -   -   -   -   -   -   -   -   -   -   -   -   -   -   -   -   -   -   -   -   -   -   -   -   -   -   -   -   -   -   -   -   -   -  11
38       -   -   -   -   -   -   -   -   -   -   -   -   -   -   -   -   -   -   -   -   -   -   -   -   -   -   -   -   -   -   -   -   -   -   -   -   -   -   -   -   -   -   -   -   -   -   -   -   -   -   -   -   -   -   -   -   -   -   -   -   -   -   -   -   -   -   -   -   -   -   -   -   -  12
39       -   -   -   -   -   -   -   -   -   -   -   -   -   -   -   -   -   -   -   -   -   -   -   -   -   -   -   -   -   -   -   -   -   -   -   -   -   -   -   -   -   -   -   -   -   -   -   -   -   -   -   -   -   -   -   -   -   -   -   -   -   -   -   -   -   -   -   -   -   -   -   -   -  13
40       -   -   -   -   -   -   -   -   -   -   -   -   -   -   -   -   -   -   -   -   -   -   -   -   -   -   -   -   -   -   -   -   -   -   -   -   -   -   -   -   -   -   -   -   -   -   -   -   -   -   -   -   -   -   -   -   -   -   -   -   -   -   -   -   -   -   -   -   -   -   -   -   -  14
41       -   -   -   -   -   -   -   -   -   -   -   -   -   -   -   -   -   -   -   -   -   -   -   -   -   -   -   -   -   -   -   -   -   -   -   -   -   -   -   -   -   -   -   -   -   -   -   -   -   -   -   -   -   -   -   -   -   -   -   -   -   -   -   -   -   -   -   -   -   -   -   -   -  15
42       -   -   -   -   -   -   -   -   -   -   -   -   -   -   -   -   -   -   -   -   -   -   -   -   -   -   -   -   -   -   -   -   -   -   -   -   -   -   -   -   -   -   -   -   -   -   -   -   -   -   -   -   -   -   -   -   -   -   -   -   -   -   -   -   -   -   -   -   -   -   -   -   -  16
43       -   -   -   -   -   -   -   -   -   -   -   -   -   -   -   -   -   -   -   -   -   -   -   -   -   -   -   -   -   -   -   -   -   -   -   -   -   -   -   -   -   -   -   -   -   -   -   -   -   -   -   -   -   -   -   -   -   -   -   -   -   -   -   -   -   -   -   -   -   -   -   -   -  17
\end{verbatim}
\end{landscape}

\section{Resultant Token Stream for the Example Test Program}
The example test program
(\texttt{src/lexer/test\_files/phase1\_program.java}) is:

\begin{lstlisting}[basicstyle=\ttfamily\small]
int sum , count , pass , mnt; while (pass !=10)
{
pass = pass + 1 ;
}
\end{lstlisting}

Running the lexer on it produces the following stream of token classes, one per
line (verbatim program output):

\lstinputlisting[basicstyle=\ttfamily\small]{example_tokens.txt}

Reading the stream against the source text, each lexeme classifies as:
\begin{center}
\small
\begin{tabular}{ll@{\hspace{3em}}ll}
\toprule
Lexeme & Token & Lexeme & Token \\
\midrule
\texttt{int}   & \texttt{int} (keyword) & \texttt{!=}   & \texttt{relop} \\
\texttt{sum}   & \texttt{id}            & \texttt{10}   & \texttt{num} \\
\texttt{,}     & \texttt{,}             & \texttt{)}    & \texttt{)} \\
\texttt{count} & \texttt{id}            & \texttt{\{}   & \texttt{\{} \\
\texttt{pass}  & \texttt{id}            & \texttt{pass} & \texttt{id} \\
\texttt{mnt}   & \texttt{id}            & \texttt{=}    & \texttt{assign} \\
\texttt{;}     & \texttt{;}             & \texttt{+}    & \texttt{addop} \\
\texttt{while} & \texttt{while} (keyword) & \texttt{1}  & \texttt{num} \\
\texttt{(}     & \texttt{(}             & \texttt{;} \texttt{\}} & \texttt{;} \texttt{\}} \\
\bottomrule
\end{tabular}
\end{center}

Note the interesting cases: \texttt{while} and \texttt{int} are reported as
reserved words, not identifiers (priority resolution inside the DFA);
\texttt{!=10} is correctly split into \texttt{relop} and \texttt{num}
(longest match); and \texttt{pass} stays an \texttt{id} even though it shares a
prefix path with no keyword. The output is byte-identical to the expected file
(\texttt{diff -u src/lexer/test\_files/phase1\_expected.tokens build/phase1.tokens}
reports no differences). The printed stream contains only the token
\emph{class} names, as the assignment requires; internally every
\texttt{ScannedToken} also carries the lexeme, its byte offset, and (for
identifiers) its symbol-table index, which the parser phase can consume.

\section{Assumptions Made and Their Justification}
\begin{enumerate}
    \item \textbf{Keywords are reserved and beat identifiers.} Keyword rules get
          priorities $0..k-1$ before all other rules, and the DFA builder picks
          the smallest priority when a state accepts several rules.
          \emph{Justification:} the assignment requires \texttt{if},
          \texttt{while}, \texttt{int}, etc.\ to be recognized as reserved words;
          resolving this once in the automaton is cleaner and faster than
          post-filtering identifier lexemes in the scanner.
    \item \textbf{Ties between equal priorities go to the earlier rule in the
          file.} \emph{Justification:} this mirrors the ``first rule wins''
          convention of Lex/Flex and makes the lexer's behavior deterministic
          and predictable from the rule file alone.
    \item \textbf{Missing DFA transitions represent the dead state.} No explicit
          error state is stored; the table prints `\texttt{-}'.
          \emph{Justification:} the dead state and all edges into/out of it carry
          no information; omitting it keeps the table sparse and readable and is
          a standard compact automaton representation. The minimizer treats a
          missing edge as target block $-1$, so correctness is preserved.
    \item \textbf{The scanner uses maximal munch (longest match).}
          \emph{Justification:} this is the universal lexical convention; without
          it, \texttt{<=} would wrongly split into \texttt{<} and \texttt{=},
          and \texttt{10} into \texttt{1} and \texttt{0}. The push-back deque
          makes the required backtracking exact and cheap.
    \item \textbf{No token may match the empty string; \texttt{\textbackslash L}
          is only valid as an alternative inside a longer pattern} (e.g.\ the
          optional exponent in \texttt{num}). The NFA builder rejects
          empty-matching rules and the scanner refuses a DFA whose start state
          accepts. \emph{Justification:} an empty match would let the scanner
          emit infinitely many tokens without consuming input.
    \item \textbf{Whitespace (space, tab, CR, LF, form feed, vertical tab)
          separates tokens and is otherwise ignored; inside expressions in the
          rule file, unescaped whitespace is formatting only.}
          \emph{Justification:} this matches both the behavior expected from the
          example program and the handout's rule-file notation, where
          \texttt{digit+ . digits} denotes concatenation of three units.
    \item \textbf{Lexical errors are recovered in panic mode by discarding one
          byte and continuing; offsets are reported on \texttt{stderr}.}
          \emph{Justification:} a scanner should find \emph{all} lexical errors in
          one pass rather than stopping at the first; discarding exactly one
          character guarantees progress, and the remaining input is rescanned so
          no valid token after the error is lost.
    \item \textbf{Single-character words in rule expressions are literals unless
          they name a definition, and range endpoints are always literal; bare
          words of two or more characters must name a definition.}
          \emph{Justification:} this resolves the ambiguity of the handout
          notation (\texttt{a-z} must mean the character range even though
          \texttt{a} could be a name) while still catching typos --
          an undefined multi-letter word is an error with a line number, with
          escaped or space-separated characters and keyword groups available for
          writing literal words.
    \item \textbf{Input is treated as single-byte characters (ASCII).}
          \emph{Justification:} the Java subset in the handout only uses ASCII
          letters, digits, and operators; bytes are read unsigned so any other
          byte simply triggers the normal error recovery.
    \item \textbf{Identifiers are interned into a simple symbol table (a vector
          of distinct lexemes) during scanning.} \emph{Justification:} later
          compiler phases need symbol identity, not just the string; building it
          in the scanner costs almost nothing and keeps token objects small.
\end{enumerate}

\section{Testing and Verification}
The implementation is covered by unit and end-to-end tests under
\texttt{src/lexer/tests/}, all of which pass:
\begin{itemize}
    \item \texttt{test\_nfa}: regex parsing and Thompson construction, including
          error cases (unmatched parentheses, dangling operators, empty
          alternatives, empty-matching rules, invalid ranges).
    \item \texttt{test\_dfa}: $\varepsilon$-closure, subset construction,
          priority resolution, minimization correctness, and table output.
    \item \texttt{test\_scanner}: longest match, push-back, symbol table, error
          offsets, and DFA validation.
    \item \texttt{test\_e2e}: the full pipeline on the Phase~1 rules and test
          program, compared against the expected token stream.
    \item \texttt{test\_flex\_parity}: the Flex-generated scanner (bonus) produces
          the same token stream and error count as the hand-written one.
\end{itemize}

\section{Bonus Task: Flex/Lex-Based Generator}
The bonus task asks how a lexical analyzer can be generated automatically with
Lex/Flex from the same set of regular expressions. Flex accepts a
pattern--action specification and generates the whole DFA-based scanner,
automating exactly the pipeline implemented by hand in this project.

\subsection{Required steps}
\begin{enumerate}
    \item Express every token class (keywords, identifiers, numbers, operators,
          punctuation) as a Flex regular expression, reusing named definitions
          (\texttt{letter}, \texttt{digit}, \texttt{digits}) that mirror the rule
          file.
    \item Place the patterns in \texttt{bonus\_flex/lexer.l}, listing keyword
          patterns \emph{before} the identifier pattern -- Flex breaks
          equal-length matches by rule order, which reproduces our keyword
          priority scheme.
    \item Attach an action to each pattern that prints the token-class name
          (or the lexeme itself for keywords and punctuation, matching the
          expected output format).
    \item Add a whitespace rule with an empty action and a catch-all \texttt{.}
          rule that counts lexical errors, mirroring the hand-written scanner's
          panic mode.
    \item Run \texttt{flex} on the specification and compile the generated
          \texttt{lex.yy.c} together with a small driver.
    \item Run the generated scanner on the example program and compare the token
          stream with the hand-written lexer (automated by
          \texttt{test\_flex\_parity}).
\end{enumerate}

\subsection{The Flex specification used}
\begin{lstlisting}[basicstyle=\ttfamily\small]
%option noyywrap

%{
#include <stdio.h>
int flex_error_count = 0;
%}

letter      [a-zA-Z]
digit       [0-9]
digits      {digit}+

%%
"boolean"|"int"|"float"|"if"|"else"|"while"     { fprintf(yyout, "%s\n", yytext); }
{letter}({letter}|{digit})*                     { fprintf(yyout, "id\n"); }
{digit}+|{digit}+\.{digits}(E{digits})?         { fprintf(yyout, "num\n"); }
"=="|"!="|">"|">="|"<"|"<="                     { fprintf(yyout, "relop\n"); }
"="                                             { fprintf(yyout, "assign\n"); }
"+"|"-"                                         { fprintf(yyout, "addop\n"); }
"*"|"/"                                         { fprintf(yyout, "mulop\n"); }
";"|","|"("|")"|"{"|"}"                         { fprintf(yyout, "%s\n", yytext); }
[ \t\r\n\f\v]+                                  { /* skip whitespace */ }
.                                               { ++flex_error_count; }
%%
\end{lstlisting}

Build and run:
\begin{lstlisting}[basicstyle=\ttfamily\small]
flex bonus_flex/lexer.l
cc -o flex_scanner lex.yy.c bonus_flex/flex_main.c
./flex_scanner src/lexer/test_files/phase1_program.java
\end{lstlisting}

\subsection{Comparison with the hand-written lexer}
Both scanners implement the same theory -- regular expressions compiled to a
minimal DFA driven with maximal munch -- and produce identical output on the
test program. The differences are practical: Flex writes the automaton for us
(less effort, fewer bugs, easy to extend with new tokens), generates highly
optimized table-compressed C code, and encodes priority by rule order; the
hand-written version, on the other hand, accepts the assignment's own rule-file
notation at \emph{run time} (no recompilation when rules change), exposes the
intermediate NFA/DFA so the transition table can be printed, and gives full
control over error reporting and the token API. Writing the lexer by hand is
what makes the generated automaton inspectable and explainable, which is the
point of this assignment, while the Flex version demonstrates how the same
specification is handled by an industrial generator.

\section{Conclusion}
The project delivers a complete, tested scanner generator and scanner for the
Java subset: rules are parsed and validated with precise error messages,
compiled through Thompson's construction and subset construction into a DFA,
minimized to 44 states with partition refinement, and driven by a
maximal-munch scanner with panic-mode recovery, a symbol table, and byte-exact
error offsets. The generated transition table and the token stream for the
example program both match the assignment's expected results, and the bonus
Flex implementation reproduces the same behavior automatically, confirming the
correctness of the hand-written pipeline.

\end{document}
